% Algebra 2 Lesson 4-4 fast-track packet: Adding and Subtracting Rational Expressions
% (enVision A2 2024 TE, Example 3 and Example 6 only; re-cut 2026-10-31 toward the Topic 4 test:
%  Topic Assessment item 10 is Example 3, item 15 is Example 6, item 4 is Example 4 (cut).
%  Practice block under R5 (quiz item 5 and topic item 15 substitute into a reciprocal-sum formula,
%  which Try It 6 does not rehearse). DOK 3 = Practice & Problem Solving 14 (Example 6), the Exit question.
%  PLC has not met on this lesson.)
% Shared preamble for the "fast" (hacked-down) Algebra 2 lesson packets, Oct 2026.
% Each packet: % Shared preamble for the "fast" (hacked-down) Algebra 2 lesson packets, Oct 2026.
% Each packet: % Shared preamble for the "fast" (hacked-down) Algebra 2 lesson packets, Oct 2026.
% Each packet: \input{fast-packet-template} then \begin{document} ... \end{document}.
\documentclass[12pt]{article}
\usepackage[margin=0.7in]{geometry}
\usepackage{amsmath,amssymb}
\usepackage{booktabs,tabularx,array}
\usepackage{enumitem}
\usepackage[T1]{fontenc}
\usepackage[utf8]{inputenc}
\usepackage{lmodern}
\usepackage{tikz}
\usepackage{pgfplots}
\pgfplotsset{compat=1.18}
\usepackage{xcolor}
\setlist{nosep,leftmargin=1.4em}
\setlength{\parindent}{0pt}
\setlength{\parskip}{6pt}
\renewcommand{\arraystretch}{1.35}
\newcolumntype{Y}{>{\raggedright\arraybackslash}X}
\pagestyle{empty}

% blanks and workspace
\newcommand{\blank}[1][2.6cm]{\rule[-2pt]{#1}{0.5pt}}
\newcommand{\lines}[1][1]{\par\vspace{2pt}\foreach \i in {1,...,#1}{\rule{0pt}{22pt}\rule{\linewidth}{0.4pt}\par}}
\newcommand{\wspart}[2]{\vspace{4pt}{\large\bfseries #1: #2}\par\vspace{-6pt}\rule{\textwidth}{0.4pt}\par}

% header: \packethead{Lesson 1-3}{Piecewise-Defined Functions}
\newcommand{\packethead}[2]{%
{\LARGE\bfseries #2}\hfill Name \blank[5cm]\\[2pt]
Algebra 2 \quad #1 (fast track) \hfill Date \blank[2.5cm] \quad Period \blank[1.2cm]\par\vspace{2pt}}

% boxed goals: \goals{learning goal}{language goal}
\newcommand{\goals}[2]{%
\fbox{\parbox{\dimexpr\textwidth-2\fboxsep-2\fboxrule}{%
\textbf{Learning goal:} #1\\[2pt]
\textbf{Language goal:} #2}}\par}

% quiz target box: \quiztarget{text}
\newcommand{\quiztarget}[1]{%
\vspace{4pt}\colorbox{gray!12}{\parbox{\dimexpr\textwidth-2\fboxsep}{\textbf{On the quiz:} #1}}\par}

% IXL table row helper: code & skill & which quiz question it feeds
\newcommand{\ixlrow}[3]{\texttt{#1} & #2 & #3 \\}

% one grid style for every graph; override xmin/xmax/... per plot
\pgfplotsset{sheetgrid/.style={
  axis lines=middle, grid=both, grid style={gray!45},
  xmin=-10.5, xmax=10.5, ymin=-10.5, ymax=10.5,
  xtick={-10,-8,...,10}, ytick={-10,-8,...,10},
  x=0.28cm, y=0.28cm,
  ticklabel style={font=\scriptsize, fill=white, inner sep=1pt},
  xlabel={$x$}, ylabel={$y$},
  axis line style={thick}, clip=true}}
% unit grid for small windows
\pgfplotsset{unitgrid/.style={sheetgrid,
  xmin=-5.5, xmax=5.5, ymin=-5.5, ymax=5.5,
  xtick={-5,...,5}, ytick={-5,...,5}, x=0.5cm, y=0.5cm}}
% open and closed endpoints
\tikzset{closed/.style={circle, fill, inner sep=1.6pt},
         open/.style={circle, draw, fill=white, inner sep=1.4pt}}
 then \begin{document} ... \end{document}.
\documentclass[12pt]{article}
\usepackage[margin=0.7in]{geometry}
\usepackage{amsmath,amssymb}
\usepackage{booktabs,tabularx,array}
\usepackage{enumitem}
\usepackage[T1]{fontenc}
\usepackage[utf8]{inputenc}
\usepackage{lmodern}
\usepackage{tikz}
\usepackage{pgfplots}
\pgfplotsset{compat=1.18}
\usepackage{xcolor}
\setlist{nosep,leftmargin=1.4em}
\setlength{\parindent}{0pt}
\setlength{\parskip}{6pt}
\renewcommand{\arraystretch}{1.35}
\newcolumntype{Y}{>{\raggedright\arraybackslash}X}
\pagestyle{empty}

% blanks and workspace
\newcommand{\blank}[1][2.6cm]{\rule[-2pt]{#1}{0.5pt}}
\newcommand{\lines}[1][1]{\par\vspace{2pt}\foreach \i in {1,...,#1}{\rule{0pt}{22pt}\rule{\linewidth}{0.4pt}\par}}
\newcommand{\wspart}[2]{\vspace{4pt}{\large\bfseries #1: #2}\par\vspace{-6pt}\rule{\textwidth}{0.4pt}\par}

% header: \packethead{Lesson 1-3}{Piecewise-Defined Functions}
\newcommand{\packethead}[2]{%
{\LARGE\bfseries #2}\hfill Name \blank[5cm]\\[2pt]
Algebra 2 \quad #1 (fast track) \hfill Date \blank[2.5cm] \quad Period \blank[1.2cm]\par\vspace{2pt}}

% boxed goals: \goals{learning goal}{language goal}
\newcommand{\goals}[2]{%
\fbox{\parbox{\dimexpr\textwidth-2\fboxsep-2\fboxrule}{%
\textbf{Learning goal:} #1\\[2pt]
\textbf{Language goal:} #2}}\par}

% quiz target box: \quiztarget{text}
\newcommand{\quiztarget}[1]{%
\vspace{4pt}\colorbox{gray!12}{\parbox{\dimexpr\textwidth-2\fboxsep}{\textbf{On the quiz:} #1}}\par}

% IXL table row helper: code & skill & which quiz question it feeds
\newcommand{\ixlrow}[3]{\texttt{#1} & #2 & #3 \\}

% one grid style for every graph; override xmin/xmax/... per plot
\pgfplotsset{sheetgrid/.style={
  axis lines=middle, grid=both, grid style={gray!45},
  xmin=-10.5, xmax=10.5, ymin=-10.5, ymax=10.5,
  xtick={-10,-8,...,10}, ytick={-10,-8,...,10},
  x=0.28cm, y=0.28cm,
  ticklabel style={font=\scriptsize, fill=white, inner sep=1pt},
  xlabel={$x$}, ylabel={$y$},
  axis line style={thick}, clip=true}}
% unit grid for small windows
\pgfplotsset{unitgrid/.style={sheetgrid,
  xmin=-5.5, xmax=5.5, ymin=-5.5, ymax=5.5,
  xtick={-5,...,5}, ytick={-5,...,5}, x=0.5cm, y=0.5cm}}
% open and closed endpoints
\tikzset{closed/.style={circle, fill, inner sep=1.6pt},
         open/.style={circle, draw, fill=white, inner sep=1.4pt}}
 then \begin{document} ... \end{document}.
\documentclass[12pt]{article}
\usepackage[margin=0.7in]{geometry}
\usepackage{amsmath,amssymb}
\usepackage{booktabs,tabularx,array}
\usepackage{enumitem}
\usepackage[T1]{fontenc}
\usepackage[utf8]{inputenc}
\usepackage{lmodern}
\usepackage{tikz}
\usepackage{pgfplots}
\pgfplotsset{compat=1.18}
\usepackage{xcolor}
\setlist{nosep,leftmargin=1.4em}
\setlength{\parindent}{0pt}
\setlength{\parskip}{6pt}
\renewcommand{\arraystretch}{1.35}
\newcolumntype{Y}{>{\raggedright\arraybackslash}X}
\pagestyle{empty}

% blanks and workspace
\newcommand{\blank}[1][2.6cm]{\rule[-2pt]{#1}{0.5pt}}
\newcommand{\lines}[1][1]{\par\vspace{2pt}\foreach \i in {1,...,#1}{\rule{0pt}{22pt}\rule{\linewidth}{0.4pt}\par}}
\newcommand{\wspart}[2]{\vspace{4pt}{\large\bfseries #1: #2}\par\vspace{-6pt}\rule{\textwidth}{0.4pt}\par}

% header: \packethead{Lesson 1-3}{Piecewise-Defined Functions}
\newcommand{\packethead}[2]{%
{\LARGE\bfseries #2}\hfill Name \blank[5cm]\\[2pt]
Algebra 2 \quad #1 (fast track) \hfill Date \blank[2.5cm] \quad Period \blank[1.2cm]\par\vspace{2pt}}

% boxed goals: \goals{learning goal}{language goal}
\newcommand{\goals}[2]{%
\fbox{\parbox{\dimexpr\textwidth-2\fboxsep-2\fboxrule}{%
\textbf{Learning goal:} #1\\[2pt]
\textbf{Language goal:} #2}}\par}

% quiz target box: \quiztarget{text}
\newcommand{\quiztarget}[1]{%
\vspace{4pt}\colorbox{gray!12}{\parbox{\dimexpr\textwidth-2\fboxsep}{\textbf{On the quiz:} #1}}\par}

% IXL table row helper: code & skill & which quiz question it feeds
\newcommand{\ixlrow}[3]{\texttt{#1} & #2 & #3 \\}

% one grid style for every graph; override xmin/xmax/... per plot
\pgfplotsset{sheetgrid/.style={
  axis lines=middle, grid=both, grid style={gray!45},
  xmin=-10.5, xmax=10.5, ymin=-10.5, ymax=10.5,
  xtick={-10,-8,...,10}, ytick={-10,-8,...,10},
  x=0.28cm, y=0.28cm,
  ticklabel style={font=\scriptsize, fill=white, inner sep=1pt},
  xlabel={$x$}, ylabel={$y$},
  axis line style={thick}, clip=true}}
% unit grid for small windows
\pgfplotsset{unitgrid/.style={sheetgrid,
  xmin=-5.5, xmax=5.5, ymin=-5.5, ymax=5.5,
  xtick={-5,...,5}, ytick={-5,...,5}, x=0.5cm, y=0.5cm}}
% open and closed endpoints
\tikzset{closed/.style={circle, fill, inner sep=1.6pt},
         open/.style={circle, draw, fill=white, inner sep=1.4pt}}

% local header: the TE title is too long to share a line with Name (template macro unchanged)
\renewcommand{\packethead}[2]{%
{\LARGE\bfseries #2}\par\vspace{4pt}
Algebra 2 \quad #1 (fast track) \hfill Name \blank[4cm] \quad Date \blank[2cm] \quad Period \blank[0.9cm]\par\vspace{2pt}}
\begin{document}

\packethead{Lesson 4-4}{Adding and Subtracting Rational Expressions}
\goals{I can find the sum or difference of rational expressions and write a compound fraction in a simpler form.}{I can explain each step of adding, subtracting or simplifying using the words factor, least common denominator, numerator and compound fraction.}

\quiztarget{Question 1 adds with different denominators: \textbf{Example 3}. Questions 2 (an LCM) and 3 (a subtraction) use its steps (full Savvas quiz only). Questions 4 and 5 simplify a compound fraction: \textbf{Example 6}. The short form grades 1 and 4 (10 points each); 5 is a +5 bonus.}
\topictarget{4}{Question 10 asks you to add $1/(x+3)$ and $-6/(x^2-9)$ (Example 3). Question 15 asks for $1/a$ when $a=1/x+1/y$ (Example 6, like the Practice block). Question 4 (subtract over $x-3$ and $3-x$) comes from Example 4, which this packet skips.}

\wspart{Example 3}{Add Rational Expressions With Unlike Denominators}
What is the sum of $\dfrac{x+3}{x^2-1}$ and $\dfrac{2}{x^2-3x+2}$? (Use the steps for adding number fractions.)

\vspace{-12pt}
\begin{align*}
&=\frac{x+3}{(x+1)(x-1)}+\frac{2}{(x-1)(x-2)} && \text{\small\textbf{1.} Factor each denominator.}\\[1pt]
&=\frac{(x+3)(x-2)}{(x+1)(x-1)(x-2)}+\frac{2(x+1)}{(x+1)(x-1)(x-2)} && \text{\small\textbf{2.} Build the LCD.}\\[1pt]
&=\frac{(x^2+x-6)+(2x+2)}{(x+1)(x-1)(x-2)}=\frac{x^2+3x-4}{(x+1)(x-1)(x-2)} && \text{\small\textbf{3.} Add the numerators.}\\[1pt]
&=\frac{(x+4)(x-1)}{(x+1)(x-1)(x-2)}=\frac{x+4}{(x+1)(x-2)} && \text{\small\textbf{4.} Factor again; divide out $x-1$.}
\end{align*}
\vspace{-4pt}
LCD: each different factor once, to its highest power ($x-1$ is in both, so it is used once). Excluded values come from the \emph{original} denominators: \textbf{answer} $(x+4)/\big((x+1)(x-2)\big)$, $x\neq-1,\,1,\,2$.

\textbf{Say it:} First I \blank[2cm] each denominator. The least common denominator is \blank[3.6cm]. I add the \blank[2.4cm] and keep the LCD. Then I factor again and divide out \blank[1.4cm].

\wspart{Try It 3}{Find the sum}
\begin{minipage}[t]{0.48\textwidth}
\textbf{a.} $\dfrac{x+6}{x^2-4}+\dfrac{2}{x^2-5x+6}$

\vspace{4pt}
LCD $=\blank[3cm]$
\lines[3]
Answer: \blank[2.6cm] for $x\neq\blank[1.4cm]$
\end{minipage}\hfill
\begin{minipage}[t]{0.48\textwidth}
\textbf{b.} $\dfrac{2x}{3x+4}+\dfrac{4x^2-11x-12}{6x^2+5x-4}$

\vspace{4pt}
LCD $=\blank[3cm]$
\lines[3]
Answer: \blank[2.6cm] for $x\neq\blank[1.4cm]$
\end{minipage}

\newpage

\wspart{Example 6}{Simplify a Compound Fraction}
A \textbf{compound fraction} has a fraction in its numerator, its denominator, or both (\mbox{\emph{fracci\'on compuesta}}). Write a simpler form of the compound fraction
\[\dfrac{\dfrac{1}{x}+\dfrac{2}{x+1}}{\dfrac{1}{y}}.\]

\begin{enumerate}[label=\textbf{\arabic*.}, itemsep=3pt]
\item List every small denominator: $x$, \ $x+1$, \ $y$. Their LCM (the least common denominator) is $x(x+1)y$.
\item Multiply the whole numerator and the whole denominator by $\dfrac{x(x+1)y}{1}$. That is multiplying by $1$, so the value does not change.
\[
\frac{\left[\frac{1}{x}+\frac{2}{x+1}\right]\cdot x(x+1)y}{\frac{1}{y}\cdot x(x+1)y}
=\frac{\frac{1}{x}\cdot x(x+1)y+\frac{2}{x+1}\cdot x(x+1)y}{\frac{1}{y}\cdot x(x+1)y}
\]
\item Divide out each small denominator. Every small fraction is gone:
\[
=\frac{(x+1)y+2xy}{x(x+1)}
\]
\item Simplify the numerator: $(x+1)y+2xy=xy+y+2xy=3xy+y=(3x+1)y$.
\[
=\frac{(3x+1)y}{x(x+1)}
\]
\end{enumerate}
\textbf{Answer:} $\dfrac{(3x+1)y}{x(x+1)}$ when $x\neq-1,\,0$ and $y\neq0$ (no small denominator can be $0$).

\textbf{Say it:} This is a \blank[3.4cm] because its numerator has fractions in it. I multiply the top and the bottom by the \blank[4.6cm], which is \blank[2.6cm]. Then no small fractions are left, and I simplify the \blank[2.4cm].

\wspart{Try It 6}{Simplify each compound fraction}
\begin{minipage}[t]{0.48\textwidth}
\textbf{a.} $\dfrac{\dfrac{1}{x-1}}{\dfrac{x+1}{3}+\dfrac{4}{x-1}}$

\vspace{4pt}
Small denominators: \blank[3.4cm]

LCD $=\blank[3cm]$
\lines[5]
Answer: \blank[2.6cm] for $x\neq\blank[1.4cm]$
\end{minipage}\hfill
\begin{minipage}[t]{0.48\textwidth}
\textbf{b.} $\dfrac{2-\dfrac{1}{x}}{x+\dfrac{2}{x}}$

\vspace{4pt}
Small denominators: \blank[3.4cm]

LCD $=\blank[3cm]$
\lines[5]
Answer: \blank[2.6cm] for $x\neq\blank[1.4cm]$
\end{minipage}

\newpage

% R5: structural parallel of quiz item 5 (parallel-resistor formula r = 1/(1/r1 + 1/r2) with r1 = x, r2 = x+1).
% Same structure (a number over a sum of two unit fractions, substituted from a formula), different numbers and setting.
\wspart{Practice}{Plug into a formula, then simplify}
A car drives to a town at $x$ miles per hour and drives back at $x+10$ miles per hour. Its average speed for the round trip is
\[
s=\frac{2}{\frac{1}{x}+\frac{1}{x+10}}.
\]
\textbf{a.} The small denominators are \blank[3cm]. The LCD is \blank[3cm]. Simplify $s$.
\lines[3]
\textbf{b.} Use your answer: if $x=40$, what is the average speed? Is it $45$ mph? \lines[1]

\wspart{IXL}{Practice for the quiz}
\begin{tabularx}{\textwidth}{|l|>{\hsize=.7\hsize}Y|>{\hsize=1.3\hsize}Y|}
\hline
\textbf{Code} & \textbf{Skill} & \textbf{Feeds} \\ \hline
\ixlrow{FEX $\bigstar$}{Add and subtract rational expressions}{Quiz 1; Quiz 3 (full Savvas quiz only) (Example 3)}
\ixlrow{YD8 $\bigstar$}{Simplify complex fractions}{Quiz 4 and 5 (Example 6)}
\ixlrow{UB5}{Factor quadratics}{Warm-up: step 1; Quiz 2 (LCM; full Savvas quiz only)}
\ixlrow{37N}{Simplify rational expressions}{Warm-up: the last step}
\hline
\end{tabularx}

\textbf{SmartScore 70 or above on each starred ($\bigstar$) skill replaces your quiz score.}

\wspart{Exit question}{Use structure (Practice \& Problem Solving 14)}
\fbox{\parbox{\dimexpr\textwidth-2\fboxsep-2\fboxrule}{%
\begin{minipage}[c]{0.36\textwidth}
\begin{tikzpicture}
% TE figure: arrowheads both ends, origin O, ticks labeled -4,-2,2,4; the TE prints the bottom vertical label
% as a second -2 (a2_4-4_TE.tex practice 14 placeholder), corrected here to -4. Grid at every unit per spec R10.
\begin{axis}[unitgrid, x=0.4cm, y=0.4cm, axis line style={thick, stealth-stealth},
  xtick={-5,...,5}, ytick={-5,...,5},
  xticklabels={,$-4$,,$-2$,,,,$2$,,$4$,}, yticklabels={,$-4$,,$-2$,,,,$2$,,$4$,}]
\node[font=\scriptsize, fill=white, inner sep=1pt, anchor=north east] at (axis cs:0,0) {$O$};
\node[closed, red] at (axis cs:1.3,0.7) {};
\node[closed, red] at (axis cs:-0.7,-1.3) {};
\node[font=\scriptsize, fill=white, inner sep=1pt, anchor=south west] at (axis cs:1.4,0.9) {$\left(\frac{1}{a},\frac{1}{b}\right)$};
\node[font=\scriptsize, fill=white, inner sep=1pt, anchor=east] at (axis cs:-1.0,-1.4) {$\left(-\frac{1}{b},-\frac{1}{a}\right)$};
\end{axis}
\end{tikzpicture}
\end{minipage}\hfill
\begin{minipage}[c]{0.62\textwidth}
Find the slope of the line that passes through the points shown. Express it in simplest form.

\vspace{4pt}
slope $=\dfrac{y_2-y_1}{x_2-x_1}=\blank[4cm]$
\lines[3]
\end{minipage}
}}

\textbf{Before the quiz, finish these sentences.}

To add or subtract, I first \blank[2cm] each denominator and find the \blank[4.4cm].

After I add, I factor the new \blank[2.6cm] and divide out a factor it shares with the denominator.

To simplify a \blank[3.4cm], I multiply its top and bottom by the LCD of all the small fractions.

\end{document}
